\par\Needspace{9\baselineskip}\begin{center}
\rule{0.6\linewidth}{0.4pt}\par\medskip
Göttingische gelehrte Anzeigen.\quad 1814 September 26.\par\medskip
\rule{0.6\linewidth}{0.4pt}
\end{center}\pdfbookmark[1]{Göttingische gelehrte Anzeigen. 1814 September 26.}{anzeige1814}

Der königlichen Gesellschaft wurde am 16.\ September von Professor \textsc{Gauss} eine Vorlesung mit dem Titel vorgelegt:

\begin{center}
\textit{Methodus nova integralium valores per approximationem inveniendi.}
\end{center}

Unter den verschiedenen Methoden zur näherungsweisen Bestimmung von Integralen oder, wie es in der Sprache der älteren Analytiker hieß, zur näherungsweisen Quadratur krummliniger Figuren ist die auf der Interpolation beruhende \textsc{Newton}-\textsc{Cotes}sche Methode eine der brauchbarsten. \textsc{Newton} hatte eine Lösung der Aufgabe angegeben, durch eine beliebige Anzahl gegebener Punkte eine parabolische Kurve zu ziehen. Deren stets leicht ausführbare Quadratur kann dann näherungsweise an die Stelle der Quadratur der eigentlich vorgegebenen, durch jene Punkte gehenden Kurve treten, und zwar umso genauer, je mehr Punkte man verwendet. \textsc{Newton} hatte es jedoch bei dieser allgemeinen Andeutung belassen und nur gleichsam als Beispiel für vier in gleichen Abständen liegende Ordinaten $A$, $B$, $C$, $D$ den näherungsweisen Flächeninhalt zwischen der ersten und der letzten bei deren Abstand $=R$ durch $(\frac{1}{8}A+\frac{3}{8}B+\frac{3}{8}C+\frac{1}{8}D)R$ angegeben. \textsc{Cotes}, der unabhängig davon bereits im Jahr 1707 und noch vor dem Erscheinen von \textsc{Newtons} Schrift \textit{Methodus differentialis} ähnliche Untersuchungen angestellt hatte, wurde durch die elegante Form, in der \textsc{Newton} das Endergebnis des obigen Beispiels dargestellt hatte (\textit{pulcherrima et}
\SourceLeaf{212}
\textit{utilissima regula} nennt es \textsc{Cotes}), dazu bewogen, diese Regeln weiter bis zum Fall von 11 Ordinaten auszudehnen. Die gesuchte Fläche erscheint dabei stets als Produkt der Basis, also des Abstands der äußersten Ordinaten, mit der Summe der mit bestimmten Zahlenkoeffizienten multiplizierten Ordinaten. Zwei vom Anfang und Ende gleich weit entfernte Ordinaten haben dabei stets gleiche Koeffizienten. Diese Quadraturkoeffizienten gibt \textsc{Cotes} bis zum Fall von 11 Ordinaten am Schluss der Abhandlung \textit{de methodo differentiali} an, die einen Teil der \textit{Harmonia mensurarum} bildet, ohne das Verfahren, mit dem er sie berechnet hat, näher zu erläutern. Vielleicht ist es dieser bescheidenen Kürze, mit der nur das Endergebnis dargestellt wird, zuzuschreiben, dass diese schöne und zweckmäßige Methode den Analytikern weniger bekannt und von ihnen weniger benutzt zu sein scheint, als sie es verdient.

Dieser Methode liegt durchweg die Voraussetzung gleicher Abstände zwischen den Ordinaten zugrunde. Auf den ersten Blick scheint diese Voraussetzung allerdings am einfachsten und natürlichsten zu sein. Die Frage war noch nicht aufgeworfen worden, ob es dennoch vorteilhafter sein könnte, Ordinaten in ungleichen Abständen zugrunde zu legen. Zur Entscheidung dieser Frage musste zunächst die Theorie der Quadraturkoeffizienten in uneingeschränkter Allgemeinheit entwickelt und der Genauigkeitsgrad des Ergebnisses bestimmt werden. Es zeigte sich, dass die Bedingungen, von denen dieser Genauigkeitsgrad abhängt, so beschaffen sind, dass man ihn durch zweckmäßig gewählte Ordinaten in ungleichen Abständen tatsächlich verdoppeln kann. Mit einer beliebigen Anzahl passend gewählter Ordinaten reicht man also ebenso weit wie mit der doppelten Anzahl von Ordinaten in gleichen Abständen. Diese Untersuchungen bilden zusammen mit der vollständigen Theorie der zweckmäßigsten Wahl der Ordinaten, der dabei zu verwendenden Quadraturkoeffizienten und der Bestimmung des Genauigkeitsgrades dieses Verfahrens den Hauptinhalt der vorliegenden Abhandlung.

Von der kurzen Entwicklung der Theorie der \textsc{Cotes}schen Quadraturkoeffizienten, die der Verfasser voranschicken zu müssen glaubte, berühren wir hier nur dasjenige, was den Genauigkeitsgrad der damit gefundenen näherungsweisen Integrale betrifft. Vor allem muss bemerkt werden, dass die Anwendbarkeit dieser Methode ebenso wie die der Interpolation voraussetzt, dass sich die Ordinaten innerhalb des zu quadrierenden Bereichs
\SourceLeaf{213}
durch eine konvergente Reihe darstellen lassen. Es sei \(x\) die Abszisse und \(y\) die Ordinate; das Integral \(\int y\,\mathrm{d}x\) werde von \(x=g\) bis \(x=h\) verlangt. Man führe anstelle von \(x\) eine andere veränderliche Größe ein, indem man etwa \(x=g+(h-g)t\) oder auch \(x=\frac12(g+h)+\frac12(h-g)u\) setzt. Dann muss sich \(y\) durch Reihen wie
\[
\alpha+\alpha't+\alpha''tt+\alpha'''t^3+\text{usw.}
\]
oder
\[
\beta+\beta'u+\beta''uu+\beta'''u^3+\text{usw.}
\]
darstellen lassen, die konvergieren, die erste zumindest, solange \(t\), die zweite, solange \(u\) nicht größer als 1 wird. Der Kürze halber kann man daher den Koeffizienten \(\alpha'\) und \(\beta'\) die Ordnung 1, den Koeffizienten \(\alpha''\) und \(\beta''\) die Ordnung 2 und so weiter zuweisen. Unter dieser Voraussetzung wird gezeigt, dass die Fehler bei der \textsc{Cotes}schen Methode zwar mit wachsender Anzahl zugrunde gelegter Werte von \(y\) stets von höherer Ordnung werden, jedoch so, dass eine ungerade Anzahl und die nächstgrößere gerade Anzahl immer Fehler derselben Ordnung hervorbringen. Für drei Ordinaten ist der Fehler also sehr nahe \(=\frac{1}{120}(h-g)\alpha''''\), für vier Ordinaten nahe \(=\frac{1}{270}(h-g)\alpha''''\), sodann für fünf Ordinaten nahe \(=\frac{1}{2688}(h-g)\alpha^{\mathrm{VI}}\) und für sechs Ordinaten nahe \(=\frac{11}{52500}(h-g)\alpha^{\mathrm{VI}}\) und so weiter. Daraus erkennt man, dass es im Allgemeinen vorteilhaft sein wird, bei der Anwendung der \textsc{Cotes}schen Methode eine ungerade Anzahl von Ordinaten zu verwenden.

Der Verfasser geht anschließend zur allgemeinen Untersuchung über, in der die Einschränkung gleicher Abstände zwischen den Ordinaten entfällt. Sind hier \(A,A',A''\) usw. die Werte von \(y\), die entsprechenden Werte von \(t\) dagegen \(a,a',a''\) usw. oder \(b,b',b''\) usw. die entsprechenden Werte von \(u\) und beträgt ihre Anzahl \(n+1\), so hat das näherungsweise Integral wiederum die Form
\[
(h-g)(RA+R'A'+R''A''+\text{usw.})
\]
Dabei sind \(R,R',R''\) usw. Zahlenkoeffizienten, die unabhängig von der Funktion \(y\) allein durch \(a,a',a''\) usw. oder durch \(b,b',b''\) usw. bestimmt werden. Die Untersuchungen des Verfassers liefern für diese Bestimmung das folgende Ergebnis. Es sei
\[
T=(t-a)(t-a')(t-a'')\,\ldotp\,\ldotp\,\ldotp\,\ldotp
\]
\SourceLeaf{214}
Die Multiplikation dieses Polynoms in \(t\), das bis zum Grad \(n+1\) reicht, mit der unendlichen Reihe, die den Logarithmus von \(\frac{t}{t-1}\) darstellt, nämlich
\[
t^{-1}+\frac12t^{-2}+\frac13t^{-3}+\text{usw.}
\]
ergebe das Produkt \(T'+T''\). Dabei bezeichnet \(T'\) das darin enthaltene Polynom in \(t\) und \(T''\) die übrige, mit negativen Potenzen von \(t\) ins Unendliche fortlaufende Reihe. Unter dieser Voraussetzung erhält man die Quadraturkoeffizienten \(R,R',R''\) usw., indem man in \(\frac{T'\mathrm{d}t}{\mathrm{d}T}\) für \(t\) der Reihe nach die Werte \(a,a',a''\) usw. einsetzt. Auf ähnliche und noch etwas bequemere Weise leitet man diese Koeffizienten aus \(b,b',b''\) usw. ab, indem man die Funktion
\[
(u-b)(u-b')(u-b'')\,\ldotp\,\ldotp\,\ldotp\,\ldotp
\]
mit \(U\) und ihr Produkt mit der unendlichen Reihe
\[
u^{-1}+\frac13u^{-3}+\frac15u^{-5}+\text{usw.}
\]
mit \(U'+U''\) bezeichnet (sodass \(U'\) das darin enthaltene Polynom in \(u\) darstellt) und dann in \(\frac{U'\mathrm{d}u}{\mathrm{d}U}\) für \(u\) der Reihe nach die Werte \(b,b',b''\) usw. einsetzt. Anstelle der rationalen Funktionen \(\frac{T'\mathrm{d}t}{\mathrm{d}T},\frac{U'\mathrm{d}u}{\mathrm{d}U}\) lassen sich auch Polynome in \(t\) und \(u\) finden, die an ihre Stelle treten können und für deren Bestimmung der Verfasser eine allgemeine Methode entwickelt.

Der Genauigkeitsgrad der Integrationsformel hängt nun von der Beschaffenheit der Reihe \(T''\) oder \(U''\) ab. Im Allgemeinen hat der Fehler zwar die Ordnung \(n+1\); fallen jedoch einige der ersten Glieder dieser Reihen weg, so hat der Fehler eine höhere Ordnung. Beginnt also \(T''\) erst mit der Potenz \(t^{-m}\) oder \(U''\) mit der Potenz \(u^{-m}\), so hat der Fehler die Ordnung \(n+m\).

Daraus ergab sich, dass man, soweit die Werte \(a,a',a''\) usw. oder \(b,b',b''\) usw. beliebig gewählt werden können, diese so bestimmen können muss, dass die ersten \(n+1\) Glieder von \(T''\) oder \(U''\) tatsächlich wegfallen. Damit erhält der Fehler der Integrationsformel die Ordnung \(2n+2\). Die Untersuchung geht deshalb dazu über, diejenigen Funktionen \(T\) und \(U\) für jeden Wert von \(n\) zu bestimmen, die die angegebene Bedingung erfüllen. Der beschränkte Raum erlaubt
\SourceLeaf{215}
uns nicht, hier auf die Einzelheiten dieser Untersuchung einzugehen: Wir bemerken deshalb nur, dass diese Funktionen einem sehr einfachen Fortschreitungsgesetz folgen und eng mit der Entwicklung der Reihen
\[
\begin{aligned}
&t^{-1}+\tfrac{1}{2}t^{-2}+\tfrac{1}{3}t^{-3}+\tfrac{1}{4}t^{-4}+\mathit{usw.}\\
&u^{-1}+\tfrac{1}{3}u^{-3}+\tfrac{1}{5}u^{-5}+\tfrac{1}{7}u^{-7}+\mathit{usw.}
\end{aligned}
\]
in Kettenbrüche zusammenhängen, die der Verfasser in einer früheren Abhandlung [\textit{Disquisitiones generales circa seriem infinitam}
\(1+\frac{\alpha.\beta}{1.\gamma}x+\mathit{usw.}\)] angegeben hat.\,---
Die Lösung der Gleichung \(T=0\) oder \(U=0\) liefert dann offenbar die Werte von \(a,a',a''\) usw. oder \(b,b',b''\) usw. Die Werte der Quadraturkoeffizienten werden nach den allgemeinen Regeln bestimmt, die in diesem Fall noch besondere Vereinfachungen zulassen. Allerdings werden in den meisten Fällen sowohl die Werte von \(a,a',a''\) usw. als auch die Quadraturkoeffizienten irrationale Größen sein. Das ist jedoch an sich ganz gleichgültig, sobald ihre Zahlenwerte ein für alle Mal mit einem angemessenen Genauigkeitsgrad berechnet sind. Ist dies geschehen, so wird die Anwendung dieser Methode auf eine beliebige Anzahl von Ordinaten wenig oder gar nicht mehr Mühe machen als die Anwendung der \textsc{Cotes}schen Methode auf dieselbe Anzahl. Letztere müsste dagegen auf eine doppelt so große Anzahl angewandt werden, um ungefähr dieselbe Genauigkeit des Ergebnisses zu liefern wie erstere.

Um für die Anwendung dieser neuen Methode nichts zu wünschen übrig zu lassen, hat der Verfasser außerdem die Zahlenwerte von \(a,a',a''\) usw. sowie von \(R,R',R''\) usw., auf 16 Dezimalstellen berechnet, mitgeteilt, zugleich mit den \textsc{Brigg}schen Logarithmen der letzteren auf 10 Dezimalstellen, alles bis zum Fall von sieben Ordinaten. In diesem letzten Fall ergibt sich für den Fehler der Integrationsformel näherungsweise \(=\frac{1}{176679360}\alpha^{\mathrm{XIV}}\). Daraus kann man schließen, dass in den meisten praktisch vorkommenden Fällen bereits eine kleinere Anzahl ausreichen wird. Um die Anwendung der Regeln und ihre relative Genauigkeit noch deutlicher zu veranschaulichen, ist als Beispiel die Berechnung von \(\int\frac{\mathrm{d}x}{\log x}\) von \(x=100000\) bis \(x=200000\) beigefügt. Bereits bei der Verwendung von vier Werten beträgt der Fehler nur \(\frac{1}{56000000}\) des Ganzen. Bei einer größeren Anzahl geht er in den unvermeidlichen Fehlern unter, die selbst die größeren Logarithmentafeln noch übrig lassen.
\par\begin{center}\rule{0.22\linewidth}{0.3pt}\end{center}
